% Part: computability
% Chapter: computability-theory
% Section: complement-ce

\documentclass[../../../include/open-logic-section]{subfiles}

\begin{document}

\olfileid{cmp}{thy}{cmp}
\olsection{مجموعه‌های محاسبه‌پذیرِ برشمردنی تحت متمم بسته نیستند}

فرض کنید $A$ محاسبه‌پذیرِ برشمردنی باشد. آیا متمم~$A$، یعنی
$\Complement{A}=\Nat\setminus A$، نیز همواره محاسبه‌پذیرِ برشمردنی
است؟ قضیه و نتیجهٔ زیر نشان می‌دهند پاسخ «نه» است.

\begin{thm}
\ollabel{thm:ce-comp}
بگذارید $A$ هر مجموعه‌ای از اعداد طبیعی باشد. آنگاه $A$ محاسبه‌پذیر
است اگر و تنها اگر $A$ و~$\Complement{A}$ هر دو محاسبه‌پذیرِ
برشمردنی باشند.
\end{thm}

\begin{proof}
جهت پیشرو آسان است: اگر $A$ محاسبه‌پذیر باشد، $\Complement{A}$ نیز
محاسبه‌پذیر است
($\Char{A}=1\tsub\Char{\Complement{A}}$)، و پس هر دو محاسبه‌پذیرِ
برشمردنی‌اند.

در جهت دیگر، فرض کنید $A$ و~$\Complement{A}$ هر دو محاسبه‌پذیرِ
برشمردنی باشند. بگذارید $A$ دامنهٔ~$\cfind{d}$ و $\Complement{A}$
دامنهٔ~$\cfind{e}$ باشد. $h$ را چنین تعریف کنید:
\[
h(x) = \umin{s}{(T(d,x,s) \lor T(e,x,s))}.
\]
به بیان دیگر، $h$ روی ورودی~$x$ در پی محاسبهٔ متوقف‌شونده‌ای از
$\cfind{d}$ یا $\cfind{e}$ می‌گردد. اگر $x\in A$ باشد در حالت نخست
موفق می‌شود، و اگر $x\in\Complement{A}$ باشد در حالت دوم. پس $h$
تابعی تام و محاسبه‌پذیر است. اکنون برای هر~$x$، $x\in A$ اگر و تنها
اگر $T(d,x,h(x))$؛ یعنی اگر $\cfind{d}$ همان تابعی باشد که تعریف شده
است. چون $T(d,x,h(x))$ رابطه‌ای محاسبه‌پذیر است، $A$ محاسبه‌پذیر است.
\end{proof}

\begin{explain}
درک آنچه رخ می‌دهد در زبان غیررسمی محاسبات آسان‌تر است: برای
تصمیم‌گیری دربارهٔ $A$، روی ورودی $x$ در پی محاسبه‌های متوقف‌شوندهٔ
$\cfind{d}$ و~$\cfind{e}$ بگردید. یکی از آن‌ها ناگزیر متوقف می‌شود؛
اگر $\cfind{d}$ باشد $x\in A$ و در غیر این صورت
$x\in\Complement{A}$.
\end{explain}

\begin{cor}
\ollabel{cor:comp-k}
$\Complement{K_0}$ محاسبه‌پذیرِ برشمردنی نیست.
\end{cor}

\begin{proof}
می‌دانیم $K_0$ محاسبه‌پذیرِ برشمردنی است، اما محاسبه‌پذیر نیست. اگر
$\Complement{K_0}$ محاسبه‌پذیرِ برشمردنی بود، بنا بر
\olref{thm:ce-comp} مجموعهٔ $K_0$ محاسبه‌پذیر می‌شد، که با
\olref[ncp]{thm:K-0} تناقض دارد.
\end{proof}

\end{document}
